\documentclass{dragonfly-report}

% A report requires two titles. The first is the short title. This is used for the 
% footer text. The second is the full title. This goes on the front of the report. 
% You are only allowed to use two LaTeX macros in the title: \\ and \emph. If others
% are used there may be issues.
\title{Example}{An Example\\ Report about \emph{sheep}}

% Reports also have a subtitle. This goes below the title on the front page.
\subtitle{A report for testing purposes}

% Authors must be included for a report. They should be listed with their full names
% with each author separated by \and.
\author{John Smith \and Jane E. Smith}

% In this example I have used the subcaption package to create subfloat. These are
% when you have more than one picture within a single figure, and each gets it's own
% internal caption.
\usepackage{subcaption}


% Optionally you can set an image to go on the titlepage. If this is not set it will
% default to the wing image as a banner accross the top of the page. If an image is
% included it should have a resulution of >= 600dpi, and ideally should be sized at
% 210x110 mmm. Wider and taller images may not work correctly. If the image is shorter
% than written, the width should not be reduced. 
%\titleimage{pattern}

% Since the title image can be comprised of images under the creative commons licence 
% there are commands to attribute these images on the inside cover. The first argument
% is the name of an image file showing the licence being used. The second is a link 
% to the image. 
%\leftimage{by-nc-nd}{https://www.flickr.com/photos/naoc0804/13923065676/in/explore-2014-04-20}
%\middleimage{by-nc-nd}{https://www.flickr.com/photos/naoc0804/13923065676/in/explore-2014-04-20}
%\rightimage{by-nc-nd}{https://www.flickr.com/photos/naoc0804/13923065676/in/explore-2014-04-20}


% Dragonfly reports
\licence{\bydragonfly}

\usepackage{lipsum}

\addbibresource{test.bib}

\begin{document}
\maketitle
\tableofcontents

\summary%[Custom Summary]
This has some boring text in it

\section{The First}
\lipsum[1]


\begin{figure}[h]
  \includegraphics[width=40mm,height=40mm]{pattern}
  \caption{This is a caption that is longer than one line. It is improtant to test 
  this behavoir as this a common feature of the dragon fly templates.}
\end{figure}

\subsection{Details}

\lipsum[2]


\begin{figure}[h]

  \begin{subfigure}{0.45\textwidth}
  \includegraphics[width=40mm,height=40mm]{pattern}
  \caption{This is a caption that is longer than one line. It is improtant to test 
  this behavoir as this a common feature of the dragon fly templates.}
\end{subfigure}\qquad
  \begin{subfigure}{0.45\textwidth}
  \includegraphics[width=40mm,height=40mm]{pattern}
  \caption{This is a caption that is longer than one line. It is improtant to test 
  this behavoir as this a common feature of the dragon fly templates.}
\end{subfigure}
\end{figure}

\subsubsection{What's This?}
\lipsum[3] 
\paragraph{A minor heading}
\lipsum[5] 
\subparagraph{A really minor heading}
\lipsum[4]

Apa comes with a number of ways of citing things.
The interesting ones are:


\begin{itemize}
  \item \cite{abraham_summary_98-09}
  \item \citet{baker_nzclassification_2010}
  \item \citep{doc_sealion_2009}
  \item \parencite{gales_phocarctos_2008}
  \item \nptextcite{ mpi_review_2012}
  %\item \fullcite{acap_summary_2011b}
  %\item \fullcitebib{acap_summary_2011b}
  \item \citeyear{robertson_population_2011}
  %\item \footcite{acap_summary_2011b}
  \item \textcite{roe_necropsy_2007}
  \item \cite{baker_marinemammals_2010}
  \item \cite{baker_census_2010}
\end{itemize}


\clearpage
\printbibliography
\end{document}
